\chapter{Modello dati e mappa del codice}
\label{chap:oggetti}

Il Capitolo~\ref{chap:flusso-codice} descrive l'ordine delle elaborazioni. Qui
sono raccolti soltanto gli oggetti scambiati tra le fasi e i moduli nei quali
intervenire. Le procedure di diagnosi e gli strumenti accessori sono invece nel
Capitolo~\ref{chap:debug}.

\section{Rappresentazioni principali}

Durante un run lo stesso componente passa da osservazione in una fotografia a
istanza fisica riconciliata, quindi a nodo del grafo. Quando si analizza
un'anomalia conviene individuare la prima rappresentazione nella quale il dato
non corrisponde più a quello atteso.

\begin{contracttable}
\code{ObservedNode} & Dataclass in \filepath{pipeline/observed_graphs.py} &
Rappresenta una detection valida in una singola immagine. Conserva codice,
confidenza, bounding box, centro 2D e identificativo della vista. \\

\code{ObservedLocalGraph} & Dataclass in
\filepath{pipeline/observed_graphs.py} & Raggruppa gli \code{ObservedNode} di
una vista. Nella baseline corrente \code{edges} e \code{local_patterns}
restano vuoti: non contiene una topologia 2D stimata. \\
\midrule

\code{SceneNode} & Dataclass in \filepath{pipeline/node_reid.py} &
Rappresenta un componente fisico ottenuto riconciliando osservazioni di viste
diverse. Mantiene identità fusa, confidenza, osservazioni di supporto e valenza
massima. Viene serializzato in \filepath{scene_nodes.json}. \\

\code{BridgeEdge} & Dataclass in
\filepath{pipeline/direct_graph_builder.py} & Descrive una connessione
candidata con evidenze geometriche e viste di supporto. È un record distinto
dall'arco NetworkX; l'elenco dei candidati viene riutilizzato dal refinement
per collegare componenti separate. \\

\code{networkx.Graph} & Struttura in memoria & È il grafo ricostruito. I nodi
derivano dagli \code{SceneNode}; gli archi conservano evidenze, forcing ed
eventuale diametro. La serializzazione produce \filepath{fused_graph.json}. \\
\end{contracttable}

\section{Moduli da conoscere}

La tabella individua i punti di intervento della pipeline senza ripetere gli
algoritmi già descritti nel capitolo precedente.

\begin{contracttable}
\filepath{shared/detection.py} & Acquisizione & Scopre gli ID disponibili,
carica YOLO e il decoder e gestisce la cache delle detection per immagine. \\

\filepath{pipeline/orchestrator.py} & Coordinamento & Implementa
\code{run_multiview_pipeline()} e \code{run_multiview_all()}; ordina le fasi,
crea le directory e salva gli artefatti. \\
\midrule

\filepath{pipeline/observed_graphs.py} & Osservazioni & Filtra le detection,
esclude i codici diametro dalla topologia e costruisce gli oggetti per vista. \\

\filepath{pipeline/node_reid.py} & ReID & Raggruppa le osservazioni in
\code{SceneNode}, costruisce la mappa osservazione--istanza e prova a
riassorbire i singleton spuri. \\

\filepath{pipeline/vggt_runtime.py},
\filepath{pipeline/vggt_export_cli.py} & Runtime VGGT & Controllano gli
artefatti disponibili e avviano, tramite l'interprete configurato, la
generazione delle point map mancanti. \\

\filepath{pipeline/vggt_3d_edges.py} & Geometria & Stima la geometria delle
osservazioni, la aggrega per \code{SceneNode} e produce gli hint usati per
valutare le connessioni. \\
\midrule

\filepath{pipeline/direct_graph_builder.py} & Grafo iniziale & Costruisce la
foresta rispettando coppie vietate, valenze e assenza di cicli; restituisce
anche i bridge candidati. \\

\filepath{pipeline/graph_refinement.py} & Completamento e output & Collega le
componenti rimaste separate, marca gli archi forzati e serializza il grafo
finale. \\

\filepath{pipeline/diameter_annotation.py} & Diametri & Associa le detection di
diametro agli archi visibili e propaga il valore lungo i rami compatibili. \\
\end{contracttable}

La configurazione operativa e i parametri globali sono riepilogati nel
Capitolo~\ref{chap:esecuzione}. Valutazione, visualizzazione e strumenti di
debug non appartengono alla pipeline principale e sono raccolti nel
Capitolo~\ref{chap:debug}.

\section{Limiti implementativi noti}

\begin{description}
  \item[Vicinato strutturale ReID.] Il vicinato viene attualmente ricavato da
  una collezione già limitata alle osservazioni dello stesso codice. Dopo
  l'esclusione del codice stesso, il vicinato risulta vuoto; la selezione delle
  fusioni dipende quindi soprattutto dal vincolo sulla stessa immagine e dalla
  distanza 3D.

  \item[Correzioni codice VGGT.] Nei rami di
  \filepath{vggt_3d_edges.py} che applicano o riproducono una correzione,
  \code{COMPONENT_INFO} viene trattato come un dizionario con campo
  \code{name}, mentre in \filepath{config.py} restituisce una tupla
  \code{(valenza, nome)}. Se una correzione cambia realmente il codice, questo
  ramo può sollevare un \code{AttributeError}.
\end{description}

Questi limiti non bloccano i casi che non attivano i relativi rami, ma devono
essere coperti da casi di regressione prima di intervenire su ReID o correzioni
VGGT.
